\hypertarget{chapter-11-mathematics-beyond-the-basics}{%
\chapter{Mathematics Beyond the Basics}\label{chapter-11-mathematics-beyond-the-basics}}

Most LaTeX users learn enough mathematics typesetting to write an equation and stop there: \texttt{\$...\$} for inline material, \texttt{\textbackslash{}begin\{equation\}} for a displayed one, a handful of symbols from \texttt{amsmath}. A book with a genuine mathematical apparatus, definitions, theorems, proofs, and a consistent notation carried across hundreds of pages, needs more structure than that, not because the underlying commands are more complicated but because the structure has to be designed rather than assembled piecemeal as each new chapter happens to need something.

\hypertarget{a-consistent-theorem-system}{%
\section{A Consistent Theorem System}\label{a-consistent-theorem-system}}

\texttt{amsthm} provides the machinery for defining theorem-like environments: \texttt{\textbackslash{}newtheorem\{theorem\}\{Theorem\}} declares an environment that numbers and labels its contents automatically, and the same mechanism extends to definitions, lemmas, propositions, corollaries, remarks, and examples, each as its own environment with its own counter. The choice that matters most, and the one worth making once at the start of a project rather than adjusting chapter by chapter, is how these counters relate to each other and to the chapter structure.

Sharing a single counter across theorems, definitions, and lemmas, so that ``Definition 3.4'' is immediately followed by ``Theorem 3.5'' rather than each type of statement keeping its own independent count, makes it possible for a reader to locate any numbered statement by its position in the document without needing to know which type of statement it was. This is achieved by declaring the later environments to share the first one's counter, \texttt{\textbackslash{}newtheorem\{lemma\}{[}theorem{]}\{Lemma\}}, rather than giving each its own. For a project as large as a full book, this is close to essential; a reader trying to find ``Lemma 12'' in a document where lemmas, theorems, and definitions all count independently has no way to search for it without already knowing roughly where it appears.

Resetting this shared counter at each chapter boundary, \texttt{\textbackslash{}newtheorem\{theorem\}\{Theorem\}{[}chapter{]}}, so numbering restarts as ``Theorem 3.1'' rather than continuing an unbroken count from the start of the book, keeps numbers legible and stable under revision: inserting a new theorem partway through Chapter 5 only renumbers the rest of Chapter 5, not every subsequent chapter in the book, which matters considerably for a document under active, ongoing revision where inserting material is common.

\hypertarget{custom-theorem-styles}{%
\section{Custom Theorem Styles}\label{custom-theorem-styles}}

\texttt{amsthm}'s \texttt{\textbackslash{}theoremstyle} command controls the visual presentation of a theorem-like environment independent of its numbering: \texttt{plain} for the traditional italicized body most readers associate with formal mathematics, \texttt{definition} for upright text more appropriate to a definition or example, and \texttt{remark} for a lighter, less formal presentation suited to asides and notes. A project can and generally should use more than one style, matching the visual weight of a statement's presentation to its actual role: a central theorem set in the more emphatic \texttt{plain} style, a routine definition in the calmer \texttt{definition} style, so that a reader's eye can distinguish the two categories of content before reading a single word of either.

Beyond the three built-in styles, a custom style defined with \texttt{\textbackslash{}newtheoremstyle} allows control over spacing, the font used for the statement's body, whether the heading and body run together or the heading stands on its own line, and how the theorem's number is formatted and punctuated. A boxed presentation, wrapping a theorem's content in a visible border, is not part of \texttt{amsthm} directly but is straightforwardly layered on top of it using \texttt{mdframed} or \texttt{tcolorbox} around a theorem environment, giving central results in a book real visual weight on the page, distinct from the plainer treatment appropriate to routine lemmas supporting a larger argument. The right choice between boxed and unboxed presentation is a document-design decision as much as a technical one: a book that boxes every single theorem-like statement loses the ability to use boxing to signal genuine importance, since nothing stands out if everything is emphasized equally.

\hypertarget{custom-operators-delimiters-and-notation-discipline}{%
\section{Custom Operators, Delimiters, and Notation Discipline}\label{custom-operators-delimiters-and-notation-discipline}}

Chapter 10 covered the mechanics of defining a new operator or symbol cleanly, avoiding collision with existing definitions. This section is about the design discipline that should govern which symbols get defined in the first place, for a project whose mathematics recurs across many chapters or many books.

A notation system chosen once, early, and used consistently thereafter is worth the up-front effort of deciding it carefully, because changing notation partway through a long project is expensive: every prior use has to be found and updated, and any reader already familiar with the earlier notation has to relearn the replacement. Deciding, before drafting begins in earnest, what symbol represents a recurring concept, what delimiters are used for what kind of grouping, and what conventions govern subscripts and superscripts across the whole project avoids the far more painful alternative of discovering three different ad hoc notations for the same idea scattered across different chapters, each introduced locally by whichever chapter happened to need it first without checking what the rest of the book was already doing.

\texttt{\textbackslash{}DeclarePairedDelimiter}, from \texttt{mathtools}, is worth knowing specifically for this kind of disciplined notation design: it defines a delimiter pair, such as a norm or an absolute value, as a single command that automatically produces correctly sized delimiters (\texttt{\textbackslash{}left}/\texttt{\textbackslash{}right}-scaled) without the author needing to write that scaling by hand at every use, and provides a starred variant for the cases where automatic sizing is not wanted. A project defining several custom delimiter-based notations benefits from using this consistently rather than hand-writing \texttt{\textbackslash{}left\textbackslash{}lVert\ \textbackslash{}cdot\ \textbackslash{}right\textbackslash{}rVert} at every occurrence, both for the reduced typing and, more importantly, for the guarantee that every use of that notation is formatted identically throughout the book.

\hypertarget{commutative-diagrams}{%
\section{Commutative Diagrams}\label{commutative-diagrams}}

\texttt{tikz-cd}, built on top of TikZ, is the standard tool for typesetting commutative diagrams: objects connected by labeled arrows, arranged in a grid, with support for the full range of arrow styles (injections, surjections, dashed arrows for existence claims, double-headed arrows for natural transformations) that this kind of diagram conventionally uses. For any project whose mathematics involves diagram chasing, functor squares, or any argument more naturally expressed as a diagram than as a sequence of equations, \texttt{tikz-cd} is worth learning directly rather than avoided in favor of describing the same content in prose, since a diagram genuinely does convey certain relationships, particularly the commutativity of multiple paths between the same two objects, more clearly and more compactly than an equivalent paragraph of text could.

The syntax is compact once learned: \texttt{\textbackslash{}begin\{tikzcd\}\ A\ \textbackslash{}arrow{[}r,\ "f"{]}\ \textbackslash{}arrow{[}d,\ "g"\textquotesingle{}{]}\ \&\ B\ \textbackslash{}arrow{[}d,\ "h"{]}\ \textbackslash{}\textbackslash{}\ C\ \textbackslash{}arrow{[}r,\ "k"\textquotesingle{}{]}\ \&\ D\ \textbackslash{}end\{tikzcd\}} produces a two-by-two commutative square with each arrow labeled, and extending this to larger diagrams is a matter of adding rows and columns to the same grid rather than learning new syntax. For a book that uses diagrams of this kind more than occasionally, defining a small set of project-level conventions, consistent arrow styles for consistent relationships, a consistent convention for where labels sit relative to their arrows, pays off the same way notation discipline does elsewhere in this chapter: a reader who has seen one of the book's diagrams can read the rest without relearning its visual vocabulary each time.

With a full mathematical apparatus in place, Part V turns to a different kind of scale problem: managing bibliographies, indexes, and glossaries across a project large enough that keeping any of them consistent by hand is no longer realistic.
